Let \(\GL_n \F_q\) denote the group of \(n \times n\) invertible matrices with entries in the finite field \(\F_q\) with \(q\) elements.
Given a matrix \(g \in \GL_n \F_q\), we define the \myemph{cycle type} of \(g\) to be \(\mu = (\mu_1,\ldots,\mu_\ell) \vdash n\) if the degrees of the irreducible factors of the characteristic polynomial of \(g\) are \(\mu_1,\ldots,\mu_\ell\) in weakly decreasing order, and we write \(\type (g) = \mu\).
This definition is built on work by Kung~\cite{kung} and Stong~\cite{stong} which suggests that a degree-\(m\) divisor of the characteristic polynomial of a matrix in \(\GL_n \F_q\) is the analog of a cycle of length \(m\) in a permutation in \(\FS_n\).
For each \(\mu \vdash n\), let \(\CT_\mu (q) \subset \GL_n \F_q\) denote the subset of matrices of cycle type \(\mu\).
Note that \(\{\CT_\mu (q) : \mu \vdash n\}\) is a partition of \(\GL_n \F_q\).

Recent work by Huang, Lewis, Morales, Reiner, and Stanton \cite{HLR, LM, refl_fact} investigated the \myemph{regular elliptic} elements of \(\GL_n \F_q\), which are those matrices whose characteristic polynomial is irreducible over \(\F_q\).
Their work suggests that the regular elliptic elements are analogous to the \(n\)-cycles in \(\FS_n\) from the perspective of enumerating factorizations.
This agrees with our definition of cycle type, as both the \(n\)-cycles of \(\FS_n\) and the regular elliptic elements of \(\GL_n \F_q\) have cycle type \((n)\).







We also consider the \myemph{regular semisimple} elements of \(\GL_n \F_q\), which are those matrices whose characteristic polynomial has no repeated irreducible factors.
Let \(\CT_\mu^\Box (q)\) denote the set of regular semisimple elements with cycle type \(\mu\).
Note that \(\{\CT_\mu^\Box (q) : \mu \vdash n\}\) is not a partition of \(\GL_n \F_q\) in general, as not all elements of \(\GL_n \F_q\) are regular semisimple for \(n \geq 2\).
However, as the following result implies, for large \(q\), an arbitrarily large proportion of \(\GL_n \F_q\) is regular semisimple.
Let \(z_\mu\) denote the cardinality of the centralizer of an element of \(\CC_\mu\) in \(\FS_n\).

\begin{corollary}[to~Corollary~\ref{rss_cc_and_cycle_type_sizes}]
\label{cycle_type_ratio_limit}
For all \(n \in \N\) and \(\mu \vdash n\),
\label{dense}
\[
\lim_{q \to \infty} \frac{\# \CT_\mu^\Box (q)}{\# \GL_n \F_q} = \lim_{q \to \infty} \frac{\# \CT_\mu (q)}{\# \GL_n \F_q} = \frac{1}{z_\mu}.
\]
\end{corollary}

In analogy with \eqref{g_kmu_defn}, we define for any \(\mu \vdash n, k \in \N\), and prime power \(q\),
\begin{align*}
g_{k,\mu} (q)
& =
\# \{ (t_1,\ldots,t_k) \in \CT_{(n)} (q)^k : t_1 \cdots t_k \in \CT_\mu (q) \}, \quad \text{and} \\
g_{k,\mu}^\Box (q)
& =
\# \{ (t_1,\ldots,t_k) \in \CT_{(n)} (q)^k : t_1 \cdots t_k \in \CT_\mu^\Box (q) \}.
\end{align*}
In this paper, we consider the quantities \(g_{k,\mu}(q)\) and \(g_{k,\mu}^\Box (q)\) to be \(\GL_n \F_q\)-analogues of \(g_{k,\mu}\),
and we seek simple formulas for computing them.
Note that \(g_{k,\mu} (q) = g_{k,\mu}^\Box (q)\) if all the parts of \(\mu\) are distinct.
